\section{Method}

\subsection{Overview}
The target of the overall movie dubbing task is:
\begin{equation}
     \hat{A} _{Dub} = \mathrm{Model}(A_{Ref}, T_{d}, V_{m}),
\end{equation}
where the $\hat{A} _{Dub}$ is the generated dubbing and $A_{Ref}, T_{d}, V_{m}$ are the reference audio, dubbing scripts, and the input silent movie clip, respectively.
The core of audio-visual alignment in movie dubbing lies in assigning accurate phoneme-level durations and prosody attributes based on the visual content and the corresponding dubbing script:
\begin{equation}
     \tilde{p}, \tilde{n}, \tilde{d} = \mathrm{Alignment}(T_d,V_{m}),
\end{equation}
where the $\tilde{p}$, $\tilde{n}$, and $\tilde{d}$ are predicted pitch, energy, and duration that align with the video content.

% main architecture 改掉,不要和Figure的caption一模一样,看着重复累赘
% The main architecture of the proposed InstructDubber is shown in Figure~\ref{method}.
Figure~\ref{method} illustrates the framework of our proposed instruction-based alignment approach.
We employ a pre-trained multimodal large language model to generate natural language fine-grained instructions of the character's speaking rate and emotional dynamics by taking both the movie clip and the script as input with the corresponding prompt.
Then, the instructed duration distilling module mines the duration cues from the speaking rate instructions, which are then combined with prosodic text features to predict phoneme-level durations $\tilde{d}$.
Meanwhile, the instructed emotion calibrating module fine-tunes a lightweight LLM-based analyzer to extract the dubbing emotions from emotion instructions, thereby guiding the prediction of emotion-aligned prosody $\tilde{p}$ and $\tilde{n}$.
We detail each module below.

\subsection{Instructed Duration Distilling}
To enable the MLLM to generate fine-grained instructions that capture the character’s speaking rate in the input video, we feed both the video clip, movie script, and the speaking rate prompt into the MLLM as input:
\begin{equation}
     I_{Dur} = \mathrm{MLLM}(T_d,V_{m},Prompt_{Dur}),
\end{equation}
where $I_{Dur}$ is the fine-grained speaking rate instruction of the given video, $Prompt_{Dur}$ is the speaking rate prompt.
Compared to conventional methods that extract visual features from each video frame, the fine-grained instructions generated by MLLMs are more informative, yet often contain redundant elements such as prepositions and conjunctions.
Therefore, it is essential to distill the discriminative duration cues from them for accurate duration prediction.

% 好好写slot attention怎么服务于电影配音任务
To alleviate this problem, we propose the Instructed Duration Distilling (IDD) module, which consists of multiple distilling layers that leverage the slot attention mechanism~\cite{slotattention}.
Slot attention initializes a set of learnable prototypes, referred to as element slots, which interact with sequence inputs to iteratively group information corresponding to the same underlying element. 
Within the IDD module, it is particularly well-suited for extracting discriminative duration cues from fine-grained speaking rate instructions, enabling effective alignment modeling.
% Slot attention initializes a set of learnable prototypes as element slots, and uses them to interact with sequence inputs and iteratively group information belonging to the same element.
% In the IDD module, slot attention is well-suited for capturing discriminative duration cues from the fine-grained speaking rate instructions.

Specifically, we first employ a global text embedding (GTE) module to convert the speaking rate instruction $I_{Dur}$ into text embeddings $E_{Dur}$ with position encoding:
\begin{equation}
     E_{Dur} = \mathrm{GTE}(T_{Dur})\in \mathbb{R}^{L_{Dur}\times d_{GTE}}.
\end{equation}
The input features of the distilling layers $E_{Dur}$ are first linearly projected to obtain key and value representations:
\begin{equation}
     K_{Dur} = W^KE_{Dur}, V_{Dur} = W^VE_{Dur}, 
\end{equation}
where $W^K$ and $W^V \in \mathbb{R}^{d_{GTE}} $ are learnable linear projections.
A fixed number $K$ of duration prototype slots $\{s^{(0)}_k\}^K_{k=1}$ are initialized randomly as learned duration features shared across all inputs.
For $T$ iterations (\ie, $T$ distilling layers), each slot $s_k^{(t)}$ is updated by attending to the input features. 
At each iteration $t$, the current slots are projected into queries:
\begin{equation}
     Q^{(t)}=W^QS^{(t)}, 
\end{equation}
where $S^{(t)} = [s_1^{(t)};...;s_K^{(t)}]$ and $W^Q$ is a learnable projection.
Attention weights between each slot and input token are computed using scaled dot-product attention:
\begin{equation}
     \alpha_{k,n}^{(t)}=\frac{\exp\left(\frac{q_k^{(t)}\cdot k_n}{\sqrt{d}}\right)}{\sum_{k^{\prime}=1}^K\exp\left(\frac{q_k^{(t)}\cdot k_n}{\sqrt{d}}\right)},\quad\mathrm{for~}k=1\ldots K,
\end{equation}
where $q_k^{(t)}\in \mathbb{R}^{d_{GTE}}$ is the $k$-th query vector and $k_n \in \mathbb{R}^{d_{GTE}}, n=1,\ldots,L_{Dur}$ is the $n$-th key vector.
Then, each slot $s_k^{(t)}$ receives the weight summary of the input values and updated via a GRU-based mechanism:
\begin{equation}
\begin{split}
    u_k^{(t)}&=\sum_{n=1}^N\alpha_{k,n}^{(t)}\cdot v_n,\quad v_n\in\mathbb{R}^{d_{GTE}},
    \\ s_k^{(t+1)}&=\mathrm{GRU}(u_k^{(t)},s_k^{(t)})+\mathrm{MLP}(\mathrm{LN}(s_k^{(t+1)})),
    \end{split}
\end{equation}
where $\mathrm{LN}$ denotes layer normalization, and $\mathrm{MLP}$ is a small feedforward network with non-linearity.
After $T$ iterations ($T$ distilling layers with shared weights), we obtain the final slots $S^{(T)}$ as the duration features $P_{Dur}$, which are distilled from the speaking rate instructions:
\begin{equation}
     P_{Dur} = S^{(T)} \in \mathbb{R}^{K\times d_{GTE}}.
\end{equation}

After getting the distilled duration features, following~\cite{ProDubber}, we convert the script text into phonemes and extract prosodic text features $T_p$ using a pretrained phoneme-level BERT model~\cite{pl=bert}:
\begin{equation}
\begin{split}
    T_{pho} &= \mathrm{G2P}(T_d)\in \mathbb{R}^{L_{pho}},
    \\ T_p &= \mathrm{BERT}_{pho}(T_{pho})\in \mathbb{R}^{L_{pho}\times d_{m}},
    \end{split}
\end{equation}
where the $\mathrm{G2P}$ and $\mathrm{BERT}_{pho}$ are the grapheme to phoneme transfer and phoneme-level BERT. 
We use the $T_p$  as queries, while the dimension-reduced duration features $P^{\prime}_{Dur}$ as both keys and values to a cross-attention (CA) layer. 
The fused duration representations are then fed into a duration predictor to obtain the final duration output:
\begin{equation}
\begin{split}
    F_{Dur}&=\mathrm{CA}(T_p, P_{Dur}^{\prime}, P_{Dur}^{\prime})\in\mathbb{R}^{L_{pho}\times d_{m}},
    \\ \tilde{d} &= \mathrm{DurationPrdictor}(F_{Dur})\in\mathbb{R}^{L_{pho}},
    \end{split}
\end{equation}
where the $\mathrm{DurationPrdictor}$ is a Bi-LSTM network with a prediction head following~\cite{StyleTTS2}.
The $ \tilde{d}$ is scaled according to the length of input video to ensure the consistency between total duration of video and dubbing.

\subsection{Instructed Emotion Calibrating}
First, we input the video, script, and prompt into the MLLM to obtain fine-grained emotion instructions:
\begin{equation}
     I_{Emo} = \mathrm{MLLM}(T_d,V_{m},Prompt_{Emo}).
\end{equation}
Similarly, it is essential to extract discriminative emotional variations from fine-grained emotion instructions to facilitate accurate emotion-aligned prosody prediction of each phoneme.
To this end and also to enhance the model's generalization ability in emotion analysis, we introduce a predefined set of emotions and use the elements as emotion entities.
After a comprehensive analysis of several emotion-labeled speech and dubbing datasets~\cite{V2C, esd}, we adopt the following seven emotions as the predefined set of emotion entities: happy, angry, disgust, fear, neutral, sad, and surprise.

We employ a lightweight LLM as the emotion instruction analyzer to extract the emotion entities from the emotion instructions.
Introducing an additional analyzer during training helps prevent the MLLM's knowledge from being biased by the visual styles of specific dubbing datasets, thereby preserving the model's generalization capability while also enhancing its flexibility to accommodate different MLLMs.
To ensure consistency between the emotion entities extracted from the emotion instruction and those present in the ground truth dubbing, we employ an audio large language model (Audio LLM) to extract ground truth emotion entities from the reference dubbing:
\begin{equation}
     Entity_{Emo} = \mathrm{AudioLLM}(A_{Dub},Prompt_{Entity}),
\end{equation}
where $Entity_{Emo}$ denotes the ground-truth emotion entities extracted from the dubbing audio.

\begin{table*}[!t]
  \centering

  \resizebox{1.0\linewidth}{!}
  {

    \begin{tabular}{c|cccc|cccc|cccc|}
    \hline
    % Setting & & \multicolumn{6}{c|}{Dub 1.0} & \multicolumn{3}{c}{Dub 2.0} \\ 
    % \hline
    
    Benchmark & \multicolumn{4}{c|}{V2C-Animation} & \multicolumn{4}{c|}{Chem} & \multicolumn{3}{c}{GRID} \\
    \midrule
    Methods 
    & DD  
    & \small{EMO-SIM} (\%) $\uparrow$
    & WER (\%) 
    & \small{UTMOS} $\uparrow$   
    & DD  
    & \small{EMO-SIM} (\%) $\uparrow$
    & WER (\%) 
    & \small{UTMOS} $\uparrow$   
    & DD  
    & WER (\%) 
    & \small{UTMOS} $\uparrow$     \\
    \midrule
    GT  & 0.0000 & 100.00 & 25.55 & 2.26 & 0.0000 & 100.00 & 3.85 & 4.18 & 0.0000 & 22.41 & 3.94\\
    % GT Mel + Vocoder & 100.00 & 25.55  & 2.26 & 99.96 & 0.00 & 0.00 & 100.00 & 22.55 & 2.26\\
    \midrule
    % V2C-Net~\cite{V2C}  & 0.5270 & 52.26 & 73.08 & 1.32 & 0.5038 & 65.84 & 90.47 & 1.81 & 0.2796 & 47.82 & 2.41 \\
    % HPMDubbing~\cite{HPMDubbing}  & 0.5690 & 49.45 & 164.16 & 1.31 & 0.6388 & 61.24 & 16.05 & 2.16 & 0.2646 & 25.51 & 2.14\\
    Speak2Dub~\cite{spk2dub} & 0.5173 & 66.58 & 17.51 & 2.41 & 0.4786 & 76.78 & 11.82 & 3.72 & 0.2650 & \textbf{17.40} & 3.69 \\
    StyleDubber~\cite{styledubber} & \textbf{0.5092} & \underline{67.22} & 31.94 & 1.89 & \underline{0.4508} & \underline{77.99} & 13.14 & 3.02 & \textbf{0.2453} & 18.88 & 3.73 \\
    DeepDubber*~\cite{deepdubber} & 0.5756 & 56.42 & 35.88 & 2.03 & 0.5041 & 52.37 & 25.51 & 2.53 & 0.3995 & 51.16 & 2.31 \\
    ProDubber~\cite{ProDubber} & 0.5148 & 67.15 & \textbf{8.04} & \underline{3.10} & 0.4673 & 76.69 & \underline{9.45} & \underline{3.85} & 0.2551 & 18.60 & \underline{3.87} \\
    \midrule
    InstructDubber (Ours) & \underline{0.5122} & \textbf{68.46} & \underline{9.27} & \textbf{3.11} & \textbf{0.4461} & \textbf{78.38} & \textbf{8.86} & \textbf{3.87} & \underline{0.2522} & \underline{17.81} & \textbf{3.88} \\
    \bottomrule
    \end{tabular}
    }
      % \vspace{-0.2cm}
        \caption{
  % \color{blue}
  Results of in-domain dubbing on three major benchmarks, which use the train set and test set from the same benchmark for training and evaluation. The best results are \textbf{in bold} and the second-best ones are \underline{underlined}.
  }

  % $\uparrow(\downarrow)$ denotes the higher (lower) value is better.
  % We follow the official code of pre-train HiFi-GAN and the extraction method of mel-spectrogram.
  \label{Ori_result}%
  % \vspace{-0.2cm}
\end{table*}%

Supervised by ground truth emotion entities, we fine-tune this analyzer using Low Rank Adaptation (LoRA)~\cite{lora} to calibrate its analysis of the emotion instruction.
For each QKV projection layer and feed-forward layer in the emotion instruction analyzer, we perform fine-tuning using rank-$R$ adaptation matrices:
\begin{equation}
    W^{\prime}=W+\Delta W=W+AB,
\end{equation}
where $W$ is the original parameter, $A\in\mathbb{R}^{d_{LLM}\times R}$ and $B\in\mathbb{R}^{R\times d_{LLM}}$ together constitute the complete set of trainable parameters $\theta$. 
During training, the parameters $\theta$ are optimized by minimizing the autoregressive loss:
\begin{equation}
    \theta^{\prime}\leftarrow\theta-\eta\cdot\nabla_\theta\mathcal{L}(\mathrm{Analyzer}(I_\mathrm{Emo};\theta),Entity_{Emo}),
\end{equation}
where $\eta$ is the learning rate.

% it 
After the fine-tuning, we use the emotion instruction analyzer to extract emotion entities from the emotion instructions and obtain their corresponding text embeddings:
\begin{equation}
     E_{Emo} = \mathrm{GTE}(\mathrm{Analyzer}(I_{Emo},Prompt_{A}, \theta^{\prime})),
\end{equation}
where $Prompt_{A}$ is the analysis prompt, $E_{Emo}\in\mathbb{R}^{L\times d_{GTE}}$ is the text embedding of the predicted emotion entities.
Then, we apply cross-attention between the dimension-reduced emotion features $E_{Emo}^{\prime}$ and the prosodic text features to obtain prosody features for prosody prediction:
\begin{equation}
\begin{split}
    F_{Emo}&=\mathrm{CA}(T_p, E_{Emo}^{\prime}, E_{Emo}^{\prime})\in\mathbb{R}^{L_{pho}\times d_{m}},
    \\ \tilde{p}, \tilde{n}&= \mathrm{ProsodyPrdictor}(F_{Emo})\in\mathbb{R}^{L_{pho}},
    \end{split}
\end{equation}
where the $\mathrm{ProsodyPrdictor}$ has the same architecture as the duration predictor with two prediction heads for pitch and energy, respectively.



\subsection{Audio Generation and Training Objective}
We feed the predicted duration and prosody, along with the script text and reference audio, into a pre-trained HiFi-GAN-based audio decoder~\cite{hifigan} to generate the final dubbing audio:
\begin{equation}
     \hat{A} _{Dub} = \mathrm{AudioDecoder}(T_{d}, \tilde{p}, \tilde{n}, \tilde{d}, A_{Ref}).
\end{equation}
The overall training objective of InstructDubber is:
\begin{equation}
    \mathcal{L}_{total} = \lambda _1\mathcal{L}_{Dur}+\lambda _2\mathcal{L}_{Emo},
    \label{loss}
\end{equation}
\begin{equation}
    \mathcal{L}_{Dur} = \frac{1}{L_{pho}} \sum_{i=0}^{L_{pho}-1} \left \| \widetilde d_i-d_i \right \|_1,
\end{equation}
\begin{equation}
\mathcal{L}_{Emo} = \frac{1}{L_{pho}} \sum_{i=0}^{L_{pho}-1} \left \| \widetilde p_i-p_i \right \|_1 + \left \| \widetilde n_i-n_i \right \|_1,
\end{equation}
where the weight in Equation~\eqref{loss} are set to $\lambda_1 = 2$, $\lambda_2 = 1$.
We use the ground truth emotion entities during training and the predicted during inference and evaluation.

